\documentclass[12pt]{article}
\usepackage[a4paper, margin=2.5cm]{geometry}
\usepackage[utf8]{inputenc}
\usepackage[T1]{fontenc}
\usepackage{lmodern}
\usepackage{hyperref}
\usepackage{graphicx}
\usepackage{setspace}
\onehalfspacing

\pagestyle{empty}

\begin{document}

\begin{flushright}
Esmeraldas, \today
\end{flushright}

\vspace{1cm}

\noindent
\textbf{Editor en Jefe} \\
Revista \textit{Ingenius} \\
Universidad Polit\'ecnica Salesiana \\
Cuenca, Ecuador

\vspace{0.8cm}

\noindent
\textbf{Asunto:} Env\'io de manuscrito original para evaluaci\'on

\vspace{0.5cm}

\noindent Estimado Editor:

\vspace{0.3cm}

\noindent Por medio de la presente, los autores abajo firmantes sometemos a consideraci\'on del Comit\'e Editorial de la Revista \textit{Ingenius} el manuscrito titulado:

\vspace{0.3cm}

\begin{center}
\textbf{\large ``Retrofitting IoT para un laboratorio remoto de cinem\'atica en ingenier\'ia: dise\~no, implementaci\'on y validaci\'on t\'ecnica de un prototipo para el estudio del MRUA en un entorno de aprendizaje h\'ibrido''}
\end{center}

\vspace{0.3cm}

\noindent para su posible publicaci\'on como \textbf{Art\'iculo Cient\'ifico} en la secci\'on de Ingenier\'ia El\'ectrica, Electr\'onica y Telecomunicaciones.

\vspace{0.3cm}

\noindent El manuscrito presenta el dise\~no, implementaci\'on y validaci\'on t\'ecnica de un prototipo de laboratorio remoto de cinem\'atica basado en \textit{retrofitting} IoT. La arquitectura propuesta integra un microcontrolador ESP32 como nodo de Edge computing, sensores infrarrojos, un broker MQTT en Raspberry Pi~4 y una interfaz web en Next.js/React. Se realiz\'o una validaci\'on cuantitativa mediante 70 ensayos independientes, demostrando que el sistema remoto alcanza precisi\'on estad\'isticamente equivalente al montaje presencial de referencia.

\vspace{0.3cm}

\noindent Declaramos que:

\begin{enumerate}
\item El manuscrito es \textbf{original e in\'edito}; no ha sido publicado previamente ni se encuentra en proceso de evaluaci\'on por otra revista o editorial.
\item Todos los autores han le\'ido y aprobado la versi\'on final del manuscrito y est\'an de acuerdo con su env\'io a la Revista \textit{Ingenius}.
\item Se han respetado los \textbf{principios \'eticos de investigaci\'on} y el trabajo est\'a libre de cualquier conflicto de inter\'es.
\item Los autores \textbf{ceden parcialmente los derechos patrimoniales} de la obra a la Universidad Polit\'ecnica Salesiana del Ecuador para ediciones impresas, conforme a la pol\'itica de la revista.
\item El trabajo se distribuir\'a bajo licencia \textbf{Creative Commons Atribuci\'on-NoComercial-SinDerivadas 4.0}.
\item Los autores son responsables del contenido y han contribuido a la concepci\'on, dise\~no, ejecuci\'on, an\'alisis e interpretaci\'on de los datos, as\'i como a la redacci\'on y revisi\'on del texto.
\end{enumerate}

\vspace{0.3cm}

\noindent \textbf{Datos de los autores:}

\vspace{0.2cm}

\noindent
\begin{tabular}{ll}
\textbf{Autor 1:} & Walter Santiago Sosa Mej\'ia \\
\textbf{ORCID:} & \href{https://orcid.org/0009-0003-6240-3462}{0009-0003-6240-3462} \\
\textbf{Afiliaci\'on:} & Pontificia Universidad Cat\'olica del Ecuador, Sede Esmeraldas \\
\textbf{Correo:} & wssosa@pucese.edu.ec \\
\\
\textbf{Autor 2 (correspondencia):} & Manuel Rogelio Nev\'arez Toledo \\
\textbf{ORCID:} & \href{https://orcid.org/0000-0001-5628-3351}{0000-0001-5628-3351} \\
\textbf{Afiliaci\'on:} & Pontificia Universidad Cat\'olica del Ecuador, Sede Esmeraldas \\
\textbf{Correo:} & manuel.nevarez@pucese.edu.ec \\
\end{tabular}

\vspace{0.8cm}

\noindent Agradecemos de antemano su consideraci\'on y quedamos a disposici\'on para cualquier informaci\'on adicional que el Comit\'e Editorial requiera.

\vspace{0.5cm}

\noindent Atentamente,

\vspace{0.5cm}

\noindent
\textbf{Walter Santiago Sosa Mej\'ia} \\
Pontificia Universidad Cat\'olica del Ecuador, Sede Esmeraldas \\
wssosa@pucese.edu.ec

\vspace{0.3cm}

\noindent
\textbf{Manuel Rogelio Nev\'arez Toledo} \\
Pontificia Universidad Cat\'olica del Ecuador, Sede Esmeraldas \\
manuel.nevarez@pucese.edu.ec

\end{document}
